\documentclass[a4paper]{article}

\usepackage[T1]{fontenc}
\usepackage[utf8]{inputenc}
\usepackage[english]{babel}

\usepackage[autostyle]{csquotes}     % Required by biblatex
\usepackage[backend=biber]{biblatex}
\addbibresource{bibliography.bib}

\usepackage{amsthm}
\usepackage{mathtools}
\usepackage{mathpartir}

\usepackage[colorlinks]{hyperref}


\theoremstyle{definition}
\newtheorem{definition}{Definition}


\newcommand*{\lvln}{n}
\newcommand*{\inddef}[2]{#1 \; \coloneq \; #2}
\newcommand*{\tyarr}[2]{#1 \rightarrow #2}
\newcommand*{\typrp}{o}
\newcommand*{\tyind}{i}
\newcommand*{\tmapp}[2]{#1 \, #2}
\newcommand*{\tmabs}[2]{\lambda \colon #1, \, #2}
\newcommand*{\tmimp}[2]{#1 \supset #2}

\newcommand*{\hastype}[3]{#1 \vdash #2 : #3}
\newcommand*{\nth}[2]{#1_{#2}}
\newcommand*{\judgement}[2]{#1 \triangleright #2}
\newcommand*{\sequent}[3]{#1; #2 \longrightarrow #3}
\newcommand*{\infcondr}[1]{\quad \text{\small #1}}
\newcommand*{\inflab}[3]{\inferrule{#1}{#2} \text{ #3}}

\newcommand*{\listappel}[2]{#1, #2}
\newcommand*{\listapp}[2]{#1 \, #2}
\newcommand*{\listcons}[2]{#1, #2}
\newcommand*{\listdelim}[1]{[#1]}
\newcommand*{\listempty}{[\;]}
\newcommand*{\listsep}{; \,}


\begin{document}


\section{Simply Typed Lambda Calculus}


\subsection{Syntax}

Taking inspiration from \cite{church:types} and \cite{pierce:plf} we define the types of the simply typed lambda calculus as:
\[
    \inddef{T}{\typrp \mid \tyind \mid \tyarr{T}{T}}
\]
That is, $\typrp$ is a type, $\tyind$ is a type and given two types $A$ and $B$, the expression $A \rightarrow B$ is also a type.
The type $\typrp$ represents the type of propositions, $\tyind$ represents the type of individuals and $A \rightarrow B$ represents the type of functions from $A$ to $B$.

Some examples:
\[
    \typrp \qquad
    \tyarr{\typrp}{\tyind} \qquad
    \tyarr{\typrp}{\tyarr{\tyind}{\typrp}} \qquad
    \tyarr{(\tyarr{\typrp}{\tyind})}{\typrp}
\]
We consider $\tyarr{}{}$ to be right associative, and thus $\tyarr{\typrp}{\tyarr{\tyind}{\typrp}}$ is the same as $\tyarr{\typrp}{(\tyarr{\tyind}{\typrp})}$.

We define the terms of the simply typed lambda calculus (using the Bruijn levels) as:
\[
    \inddef{t}{\lvln \mid \tmapp{t}{t} \mid \tmabs{T}{t} \mid \bot \mid \top \mid t \land t \mid t \lor t \mid \tmimp{t}{t}}
\]
where $\lvln$ is a natural number and $T$ is a type.

Some examples:
\[
    \tmabs{\tyind}{0} \qquad
    \tmapp{(\tmabs{\tyarr{\tyind}{\typrp}}{2})}{1} \qquad
    \tmabs{\typrp}{\tmapp{3}{\bot}} \qquad
    \tmabs{\tyarr{\typrp}{\typrp}}{\tmabs{\tyind}{\tmapp{0}{\top}}}
\]


\subsection{Typing}

Given a list of types $\Gamma$, a natural number $n$ and a type $T$, we denote the $n$-th element of $\Gamma$ by $\nth{\Gamma}{n}$ and the list constructed by appending the element $T$ to the end of $\Gamma$ as $\listappel{\Gamma}{T}$.

\begin{definition}
    Given a list of types $\Gamma$, a term $t$ and a type $T$, we define the proposition $\hastype{\Gamma}{t}{T}$ by the following inference rules:
    \begin{gather*}
        \inferrule{ }{\hastype{\Gamma}{\lvln}{T}} \infcondr{if $\nth{\Gamma}{\lvln} = T$} \\
        %
        \inferrule{\hastype{\Gamma}{s}{\tyarr{A}{B}} \\ \hastype{\Gamma}{t}{A}}{\hastype{\Gamma}{\tmapp{s}{t}}{B}} \\
        %
        \inferrule{\hastype{\listappel{\Gamma}{A}}{s}{B}}{\hastype{\Gamma}{\tmabs{A}{s}}{\tyarr{A}{B}}} \\
        %
        \inferrule{ }{\hastype{\Gamma}{\bot}{\typrp}} \qquad
        \inferrule{ }{\hastype{\Gamma}{\top}{\typrp}} \\
        %
        \inferrule{\hastype{\Gamma}{t}{\typrp} \\ \hastype{\Gamma}{s}{\typrp}}{\hastype{\Gamma}{t \land s}{\typrp}} \qquad
        \inferrule{\hastype{\Gamma}{t}{\typrp} \\ \hastype{\Gamma}{s}{\typrp}}{\hastype{\Gamma}{t \lor s}{\typrp}} \\
        %
        \inferrule{\hastype{\Gamma}{t}{\typrp} \\ \hastype{\Gamma}{s}{\typrp}}{\hastype{\Gamma}{\tmimp{t}{s}}{\typrp}}
    \end{gather*}
    If the proposition $\hastype{\Gamma}{t}{T}$ holds, we say that $t$ \emph{has type} $T$ in the \emph{context} $\Gamma$.
\end{definition}

Some examples:
\begin{gather*}
    \inferrule{ }{\hastype{\listdelim{\tyind \listsep \typrp}}{0}{\tyind}} \\
    %
    \inferrule{
        \inferrule{
            \inferrule{
                \inferrule{ }{\hastype{\listdelim{\tyarr{\tyind}{\typrp} \listsep \tyind}}{0}{\tyarr{\tyind}{\typrp}}} \\
                \inferrule{ }{\hastype{\listdelim{\tyarr{\tyind}{\typrp} \listsep \tyind}}{1}{\tyind}}
            }{
                \hastype{\listdelim{\tyarr{\tyind}{\typrp} \listsep \tyind}}{\tmapp{0}{1}}{\typrp}
            }
        }{
            \hastype{\tyarr{\tyind}{\typrp}}{\tmabs{\tyind}{\tmapp{0}{1}}}{\tyarr{\tyind}{\typrp}}
        }
    }{
        \hastype{}{\tmabs{\tyarr{\tyind}{\typrp}}{\tmabs{\tyind}{\tmapp{0}{1}}}}{\tyarr{(\tyarr{\tyind}{\typrp})}{\tyarr{\tyind}{\typrp}}}
    }
\end{gather*}

Alternatively, we could have defined $\land$ and $\lor$ as terms of type $\tyarr{\typrp}{\tyarr{\typrp}{\typrp}}$.


\section{G Logic}

Given a list of types $\sigma$ and a term $t$, the expression $\judgement{\sigma}{t}$ is called a \emph{judgement}.
Let $\Sigma$ be a list of types, $\Gamma$ a list of judgements and $B$ a judgement. The expression $\sequent{\Sigma}{\Gamma}{B}$ is called a \emph{sequent}.
We denote the list constructed by appending a head element $B$ to a tail $\Gamma$ as $\listcons{B}{\Gamma}$.
The concatenation of two lists $\Sigma$ and $\sigma$ is denoted by $\listapp{\Sigma}{\sigma}$.

We define a calculus with the following inference rules:
\begin{gather*}
    \inflab{
        \hastype{\listapp{\Sigma}{\sigma}}{t}{\typrp}
    }{
        \sequent{\Sigma}{\listdelim{\judgement{\sigma}{t}}}{\judgement{\sigma}{t}}
    }{init} \qquad
    %
    \inflab{
        \sequent{\Sigma}{\Delta}{B} \\
        \sequent{\Sigma}{\listcons{B}{\Gamma}}{C}
    }{
        \sequent{\Sigma}{\listapp{\Delta}{\Gamma}}{C}
    }{cut} \\
    %
    \inflab{
        \sequent{\Sigma}{\listcons{B}{\listcons{B}{\Gamma}}}{C}
    }{
        \sequent{\Sigma}{\listcons{B}{\Gamma}}{C}
    }{$cL$} \qquad
    %
    \inflab{
        \hastype{\listapp{\Sigma}{\sigma}}{t}{\typrp} \\
        \sequent{\Sigma}{\Gamma}{C}
    }{
        \sequent{\Sigma}{\listcons{\judgement{\sigma}{t}}{\Gamma}}{C}
    }{$wL$} \\
    %
    \inflab{
        \hastype{\listapp{\Sigma}{\sigma}}{t}{\typrp}
    }{
        \sequent{\Sigma}{\listdelim{\judgement{\rho}{\bot}}}{\judgement{\sigma}{t}}
    }{$\bot L$} \qquad
    %
    \inflab{
    }{
        \sequent{\Sigma}{\listempty}{\judgement{\sigma}{\top}}
    }{$\top R$} \\
    %
    \inflab{
        \hastype{\listapp{\Sigma}{\sigma}}{s}{\typrp} \\ \sequent{\Sigma}{\listcons{\judgement{\sigma}{t}}{\Gamma}}{B}
    }{
        \sequent{\Sigma}{\listcons{\judgement{\sigma}{t \land s}}{\Gamma}}{B}
    }{$\land L$} \qquad
    %
    \inflab{
        \hastype{\listapp{\Sigma}{\sigma}}{t}{\typrp} \\ \sequent{\Sigma}{\listcons{\judgement{\sigma}{s}}{\Gamma}}{B}
    }{
        \sequent{\Sigma}{\listcons{\judgement{\sigma}{t \land s}}{\Gamma}}{B}
    }{$\land L$} \\
    %
    \inflab{
        \sequent{\Sigma}{\Gamma}{\judgement{\sigma}{s}} \\
        \sequent{\Sigma}{\Gamma}{\judgement{\sigma}{t}}
    }{
        \sequent{\Sigma}{\Gamma}{\judgement{\sigma}{s \land t}}
    }{$\land R$} \\
    %
    \inflab{
        \sequent{\Sigma}{\listcons{\judgement{\sigma}{s}}{\Gamma}}{B} \\
        \sequent{\Sigma}{\listcons{\judgement{\sigma}{t}}{\Gamma}}{B}
    }{
        \sequent{\Sigma}{\listcons{\judgement{\sigma}{s \lor t}}{\Gamma}}{B}
    }{$\lor L$} \\
    %
    \inflab{
        \hastype{\listapp{\Sigma}{\sigma}}{t}{\typrp} \\
        \sequent{\Sigma}{\Gamma}{\judgement{\sigma}{s}}
    }{
        \sequent{\Sigma}{\Gamma}{\judgement{\sigma}{s \lor t}}
    }{$\lor R$} \qquad
    %
    \inflab{
        \hastype{\listapp{\Sigma}{\sigma}}{s}{\typrp} \\
        \sequent{\Sigma}{\Gamma}{\judgement{\sigma}{t}}
    }{
        \sequent{\Sigma}{\Gamma}{\judgement{\sigma}{s \lor t}}
    }{$\lor R$} \\
    %
    \inflab{
        \sequent{\Sigma}{\Gamma}{\judgement{\sigma}{s}} \\
        \sequent{\Sigma}{\listcons{\judgement{\sigma}{t}}{\Gamma}}{B}
    }{
        \sequent{\Sigma}{\listcons{\judgement{\sigma}{\tmimp{s}{t}}}{\Gamma}}{B}
    }{$\tmimp{}{} L$} \qquad
    %
    \inflab{
        \sequent{\Sigma}{\listcons{\judgement{\sigma}{s}}{\Gamma}}{\judgement{\sigma}{t}}
    }{
        \sequent{\Sigma}{\Gamma}{\judgement{\sigma}{\tmimp{s}{t}}}
    }{$\tmimp{}{} R$}
\end{gather*}


\printbibliography


\end{document}
